\chapter{Interventional Effects}\label{chapter:interventional_effects}

Motivated by the problems presented by situations with treatment-induced confounding,~\cite{vanderweeleetal2014} 
proposed a randomized-interventional analogue of natural effects that are identified without 
requiring the cross-world independence assumption. Instead of defining 
a unit-level direct and indirect effect estimand as in the case of natural effects, interventional effects
target population-level average outcomes from a hypothetical joint intervention on the treatment and the 
mediator. A key difference between natural and interventional effects
is that interventional effects can still be identified given a set of reasonable assumptions in cases where
there is either treatment-induced confounding or when there are multiple mediators of interest. We begin by
reviewing the definition and interpretation of interventional effects in one mediator problems before 
extending our analysis to the two mediator case. We then discuss what we
view as the practical use and interpretation of interventional effects and the difficulties with 
interpretation in multiple mediator settings specifically. In response to the issues with interpretation 
we raise, we propose an extension of the interventional effects definitions from~\cite{vanderweelerobinson2014}
as a promising avenue for future exploration of interventional effects. 

\section{One Mediator}

\subsection{Effect Definitions}\label{subsection:intv_eff_defs}

Let $G_{i}(a)$ represent a random draw of a mediator value from $M(a)$, which now represents a 
the population-level distribution of the potential mediator given treatment level $a$ (and covariates $C_{i}$).
Let $Y_{i}(a,G_{i}(a'))$ represent a potential outcome for unit $i$ assigned treatment $a$ and
a random level for the mediator from the distribution of $G(a')$. This is a stochastic potential 
outcome: the potential outcome is itself a random variable because its value depends on a realization
from the distribution of $G(a')$~\citep{vanderweelerobins2012}. 

The definitions for interventional effects mirror those for natural effects above, except that they are
defined in terms of stochastic potential outcomes $Y_{i}(a,G_{i}(a'))$, which represent joint interventions
to assign unit $i$ values for the treatment and the mediator. First, we write out the equation, the so-called 
mediational g-formula, to estimate the mean of $Y_{i}(a,G_{i}(a'))$ for any 
$a, a' \in \mathcal{T}$~\citep{kuhaetal2021}: 

\begin{equation}\label{eq:RIE_one_mediator_g_comp}
\begin{aligned}
	\mathbb{E}[Y(a,G(a'))|C] = 
	\int_{m} \mathbb{E}[Y(a,m)|T = a, M = m, C]\mathbb{P}[M(a') = m|C]dm
\end{aligned}
\end{equation}

We then define interventional effects in the following way: 

\begin{itemize}[noitemsep,topsep=0pt]
\item  $\mathbb{E}[Y(1,G(1)) - Y_{i}(0,G(0))]$ --- The \textit{interventional total effect} 
represents the effect of a joint intervention that gives a unit treatment while assigning their mediator 
level to a random draw from the distribution of the mediator for treated units, compared to an intervention
to assign a unit to no treatment while drawing their mediator level from the distribution of the mediator for 
untreated units. 
\item  $\mathbb{E}[Y(1,G(0)) - Y_{i}(0,G(0))]$ --- The \textit{pure interventional direct effect} 
represents the effect of a joint intervention that gives a unit treatment while drawing their mediator level 
from the distribution of the mediator for untreated units, compared to an intervention
to assign a unit to no treatment while drawing their mediator level from the distribution of the mediator for 
untreated units. 
\item  $\mathbb{E}[Y(1,G(1)) - Y_{i}(0,G(1))]$ --- The \textit{total interventional direct effect} 
represents the effect of a joint intervention that gives a unit treatment while drawing their mediator level 
from the distribution of the mediator for treated units, compared to an intervention
to assign a unit to no treatment while drawing their mediator level from the distribution of the mediator for 
treated units. 
\item  $\mathbb{E}[Y(0,G(1)) - Y_{i}(0,G(0))]$ --- The \textit{pure interventional indirect effect} 
represents the effect of a joint intervention that gives a unit no treatment while drawing their mediator level 
from the distribution of the mediator for treated units, compared to an intervention
to assign a unit to no treatment while drawing their mediator level from the distribution of the mediator for 
untreated units.
\item  $\mathbb{E}[Y(1,G(1)) - Y_{i}(1,G(0))]$ --- The \textit{total interventional indirect effect} 
represents the effect of a joint intervention that gives a unit treatment while drawing their mediator level 
from the distribution of the mediator for treated units, compared to an intervention
to assign a unit to treatment while drawing their mediator level from the distribution of the mediator for 
untreated units.
\end{itemize}

The total interventional effect is formally defined in~\cite{linvanderweele2017}. 
Much like with natural effects, there is a decomposition of the interventional total effect into the sum of the 
pure interventional direct effect and the total interventional indirect effect (or the total interventional
direct effect and the pure interventional indirect effect): 

\begin{equation}
\begin{aligned}
	\mathbb{E}[Y(1,G(1)] - \mathbb{E}[Y(0,G(0))] = 
	\{\mathbb{E}[Y(1,G(0))] - \mathbb{E}[Y(0,G(0))]\}\\ 
	+ \{\mathbb{E}[Y(1,G(1))] - \mathbb{E}[Y(1,G(0))]\}
\end{aligned}
\end{equation}

The pure and total interventional effects are not explicitly defined 
by~\cite{vanderweeleetal2014} and a distinction is generally not made between them in the literature, though 
the notion of a pure and total interventional effect appears in the eAppendix of~\cite{vanderweele2014}. 
The total interventional effect and the total causal effect are equivalent when the cross-world independence
assumption is satisfied, but the two quantities are not otherwise equivalent~\citep{vanderweeleetal2014}.
Conceptually, this occurs because the joint intervention on the treatment and the mediator breaks the 
causal ``link'' between $T$ and $M$ that is preserved when estimating the total causal effect.
The total interventional effect is formally defined in~\cite{linvanderweele2017}. Thus, the 
sum of a direct and indirect interventional effect does not always equal the total causal effect, which 
is sometimes viewed as a shortcoming of these effects~\citep{vansteelandtdaniel2017}. However, as we
argue in Section~\ref{subsection:interpretation}, we do not believe that such decompositions are 
particularly instructive given the interpretative goals of interventional effects. 

Additionally,~\cite{didelez2006} and ~\cite{geneletti2007} defined a form of these types of
interventional effects in terms of general randomized interventions on the mediator without defining these
effects in terms of counterfactuals. These effects are very similar to the definition of interventional effects
given here. However, they allow for any stochastic intervention on the mediator rather than
restricting such interventions to the distribution of the mediator in two treatment groups. 

The above effects are identified under consistency, composition, and conditions (i)-(iii) of 
Assumption~\ref{assm:seq_ignorability}. 
In particular, the cross-world independence assumption~\citep{vanderweeleetal2014} is not required. 

\subsection{Experimental Identification}

A key feature of interventional effects is that it is in principle possible to 
design randomized experiments that can identify these effects, and, 
given the prescriptive interpretation of interventional effects, the specific interventions carried out in such an 
experiment would be the causal quantities we are ultimately interested in. Although such experiments might be infeasible 
or unethical in practice, being able to conceptualize an ideal experiment to identify such effects can still
help inform observational study design and interpretation. 

The experimental setup, like our running example, is inspired by~\cite{morenobetancur2018}. Suppose we 
are trying to estimate interventional effects in the running example. Suppose a ``headache drug'' has been developed 
that allows us to exactly set the level of headache a patient who takes this drug will experience while having no
other side effects. We set up the following 6-arm trial, with patients randomly assigned to each arm: 

\begin{experiment}\label{exp:one_mediator}
\indent 
\begin{enumerate}[noitemsep,topsep=0pt]
	\item Patients in \textbf{Arm 1} receive the anti-depressant, while patients in \textbf{Arm 2} do not. 
	The levels of both headache and depression are recorded for each patient.
	\item After recording the results of the experiment in stage 1, patients in arms 3-6 receive the following
	combinations of treatment and headache medication: 
	\begin{enumerate}[noitemsep,topsep=0pt]
		\item \textbf{Arm 3} receives the anti-depressant, and the headache drug is used to give each patient 
		a randomly selected level of headache from the distribution of observed headache levels in Arm 1. 
		\item \textbf{Arm 4} receives the anti-depressant, and the headache drug is used to give each patient 
		a randomly selected level of headache from the distribution of observed headache levels in Arm 2. 
		\item \textbf{Arm 5} does not receive the anti-depressant, and the headache drug is used to give each patient 
		a randomly selected level of headache from the distribution of observed headache levels in Arm 1. 
		\item \textbf{Arm 6} does not receive the anti-depressant, and the headache drug is used to give each patient 
		a randomly selected level of headache from the distribution of observed headache levels in Arm 2. 
	\end{enumerate}
\end{enumerate}
\end{experiment}

\begin{table}[ht]
\centering
\caption{Table of results for Experiment~\ref{exp:one_mediator}}
\label{tab:rct_one_mediator}
\begin{tabularx}{\textwidth}{@{}
  >{\hsize=0.40\hsize\centering\arraybackslash}X
  >{\hsize=0.40\hsize\centering\arraybackslash}X
  >{\hsize=1.40\hsize\centering\arraybackslash}X
  >{\hsize=1.40\hsize\centering\arraybackslash}X
  >{\hsize=1.40\hsize\centering\arraybackslash}X @{}}
\hline
\addlinespace[6pt]
\makecell[bc]{\textbf{Stage}} &
\makecell[bc]{\textbf{Arm}} &
\makecell[bc]{\textbf{Treatment}} &
\makecell[bc]{\textbf{Headaches}} &
\makecell[bc]{\textbf{Depression}} \\
\addlinespace[3pt]
\noalign{\hrule height 1pt}
\addlinespace[3pt]
\multirow{2}{*}{1} 
&1 & 1 & $\hat{M}(1)_{1}$ & $y_{1}$\\
&2 & 0 & $\hat{M}(0)_{2}$ & $y_{2}$\\
\addlinespace[2pt]
\noalign{\hrule height 0.5pt}
\addlinespace[2pt]
\multirow{4}{*}{2}
&3 & 1 & $\hat{G}(1)_{3}$ & $y_{3}$\\
&4 & 1 & $\hat{G}(0)_{4}$ & $y_{4}$\\
&5 & 0 & $\hat{G}(1)_{5}$ & $y_{5}$\\
&6 & 0 & $\hat{G}(0)_{6}$ & $y_{6}$\\
\addlinespace[3pt]
\noalign{\hrule height 1pt}
\end{tabularx}

\vspace{4pt}
{\footnotesize
\emergencystretch=3em
\setlength{\parindent}{1.5em}
\textsuperscript{1} $\hat{M}(1)$ and $\hat{M}(0)$ represent the distributions of headache in
arms 1 and 2, respectively. $\hat{G}(1)$  and $\hat{G}(0)$ represent the distributions of assigned
headache from random draws from $\hat{M}(1)$ and $\hat{M}(0)$, respectively.

Finally, $y$ gives the observed average level of depression in each arm.\par}
\end{table}

Given this experimental setup, it is possible to identify the suite of interventional effects defined in
Section~\ref{subsection:intv_eff_defs} by taking observed differences in outcomes $y$ (see Table~\ref{tab:rct_one_mediator}). For
example, the estimated total interventional effect is $y_{3} - y_{6}$, the pure interventional direct
effect is $y_{4} - y_{6}$, and the pure interventional indirect effect is $y_{5} - y_{6}$. Such an 
experimental setup also highlights how despite the focus on the pure direct and total indirect 
interventional effects in the literature, there is no \textit{a priori} reason to prefer these 
interventional effects to the other effects that are identified by such an experiment. 

\section{Multiple Mediators}

We now extend the interventional effects framework to a situation with two mediators. Additionally, we 
assume that we know the causal ordering of these two mediators (though this condition can be relaxed in 
some cases). Extending the running example from~\cite{morenobetancur2018}, we now suppose that we notice 
that the anti-depressant, in addition to causing headache, can cause insomnia in the first few weeks of
treatment. We also believe that insomnia can cause chronic headaches, so headaches can be caused both 
directly by the anti-depressant and indirectly through the anti-depressant's effect on insomnia. The 
causal structure of this problem is represented in Figure~\ref{fig:two_mediators} (note that this is 
identical to Figure~\ref{fig:exposure_induced_confounder}). 

\begin{figure}
\centering
\begin{tikzpicture}

  % Define nodes
  \node[obs] (T) {$\mathbf{T}$};
  \node[obs, right=of T, yshift=1.5cm] (L) {$\mathbf{M_{1}}$};
  \node[obs, right=2cm of T] (M) {$\mathbf{M_{2}}$};
  \node[obs, right=1cm of M]  (Y) {$\mathbf{Y}$};
  \node[obs, left=of T]  (C) {$\mathbf{C}$};

  % Connect the nodes
  \edge {T,L} {M} ; %
  \edge {M} {Y} ; %
  \draw[->] (T) to[out=-30, in=-150] (Y); %
  \draw[->] (L) to[out=0, in=120] (Y); %
  \draw[->] (C) to[out=300, in=240] (Y); %
  \draw[->] (C) to[out=315, in=225] (M); %
  \edge {C,T} {L} ; %
  \edge {C} {T} ; %

\end{tikzpicture}
\caption{Two mediators with known causal ordering}
\label{fig:two_mediators}
\end{figure}


There are two significant definitions of interventional effects with
two mediators in the literature that we will focus on. The first, proposed by~\cite{linvanderweele2017} and
inspired by similar work done by~\cite{vanderweeletchetgentchetgen2017}, attempts to estimate 
interventional effect analogues of path-specific effects $T \rightarrow Y$, 
$T \rightarrow M_{1} \rightarrow Y$, $T \rightarrow M_{2} \rightarrow Y$ and 
$T \rightarrow M_{1} \rightarrow M_{2} \rightarrow Y$. The second proposal, by~\cite{vansteelandtdaniel2017}
has the advantage of not making any assumptions about the causal ordering of the mediators. Instead,
interventional effects are estimated from interventions on the marginal distributions of the mediators, 
assuming the mediators are independent. Afterwards, a remainder term, which represents the dependence between
the mediators, is computed. 

\subsection{Lin-VanderWeele Interventional Effects}

We now let $M_{1,i}$ be the observed level of insomnia for unit $i$ and $M_{2,i}$ the observed level of
headache. Let $\mathcal{M}_{1}$ be the support of $M_{1}$ and $\mathcal{M}_{2}$ the support of $M_{2}$. 
We also redefine potential mediators and outcomes as follows:

\begin{itemize}[noitemsep,topsep=0pt]
	\item Let $M_{1,i}(a)$ represent the potential levels of insomnia for unit $i$ given 
	treatment level $a$. 
	\item Let $M_{2,i}(a,M_{1,i}(a'))$ represent the potential levels of headache for unit $i$ given 
	treatment level $a$ and an insomnia level $M_{1,i}(a')$ that they would have if they received 
	treatment level $a'$. 
	\item Let $Y_{i}(a,M_{1,i}(a'),M_{2,i}(a'',a'''))$ represent the potential outcomes 
	for unit $i$ given treatment level $a$ and potential mediator levels $M_{1,i}(a')$ and 
	$M_{2,i}(a'',a''')$. 
\end{itemize}

Finally, we define $G_{1,i}(a)$ as a random draw from the distribution of $M_{1}(a)$ 
and\\ $G_{2,i}(a,M_{1}(a'))$ as a random draw from the distribution of $M_{2}(a,M_{1}(a'))$. 

The main idea in~\cite{linvanderweele2017} is to estimate
path-specific interventional effects in the multiple mediator setting in a similar way to how 
such effects are estimated in the single mediator case. In particular, they pursue a decomposition
of the total interventional effect in terms of path-specific interventional effects 
$T \rightarrow Y$, $T \rightarrow M_{1} \rightarrow Y$, $T \rightarrow M_{2} \rightarrow Y$, 
and $T \rightarrow M_{1} \rightarrow M_{2} \rightarrow Y$, which mirrors the decomposition of 
the total interventional effect in terms of the pure direct interventional effect and the total 
indirect interventional effect in the one mediator case. This decomposition relies on the following
simple algebraic expansion: 

\begin{equation}
\begin{aligned}
	\mathbb{E}[Y(1,G_{1}(1),G_{2}(1,G_{1}(1)))] - 
	\mathbb{E}[Y(0,G_{1}(0),G_{2}(0,G_{1}(0)))]\\ 
	= \Psi(1,1,1,1) - \Psi(0,0,0,0)\\
	= \{\Psi(1,0,0,0) - \Psi(0,0,0,0)\} + \{\Psi(1,1,0,0) - \Psi(1,0,0,0)\} +\\ 
	\{\Psi(1,1,1,0) - \Psi(1,1,0,0)\} + \{\Psi(1,1,1,1) - \Psi(1,1,1,0)\}
\end{aligned}
\end{equation}

where $\Psi$ is defined by the mediational g-formula given in~\cite{linvanderweele2017}:

\begin{equation}\label{eq:RIE_two_mediator_g_comp}
\begin{aligned}
	\Psi(a,a',a'',a''') = \mathbb{E}[Y(a,G_{1}(a'),G_{2}(a'',a'''))|C] =\\ 
	\int_{m_{2}}\int_{m_{1}'}\int_{m_{1}} \mathbb{E}[Y(a,m_{1},m_{2})|T = a, 
	M_{1} = m_{1}, M_{2} = m_{2}, C]\\
	\mathbb{P}[M_{2}(a'',m_{1}') = m_{2}|T = a''', M_{1} = m_{1}', C]\\
	\mathbb{P}[M_{1}(a''') = m_{1}|C]\mathbb{P}[M_{1}(a') = m_{1}'|C]dm_{1}dm_{1}'dm_{2}
\end{aligned}
\end{equation}

Each term in the decomposition represents a particular interventional effect that corresponds 
to the pathways 
$T \rightarrow Y$, $T \rightarrow M_{1} \rightarrow Y$, $T \rightarrow M_{2} \rightarrow Y$, 
and $T \rightarrow M_{1} \rightarrow M_{2} \rightarrow Y$, respectively. However, the mediational
g-formula allows for identification of 16 different quantities $\Psi$ for a binary treatment $a$ 
that we might be interested in. In general, we conceive of interventional 
effects with two mediators to be the average potential outcome of a joint intervention that assigns
a treatment level $a$, a level of the first mediator randomly drawn from the distribution of the
mediator in a population with treatment level $a'$, and a level of the second mediator randomly drawn
from the distribution of the mediator in a population with treatment level $a''$ and first mediator 
level randomly drawn from a population with treatment level $a'''$. 

~\cite{linvanderweele2017} proved that equation~\ref{eq:RIE_two_mediator_g_comp} identifies 
$\Psi(a,a',a'',a''')$, and thus any estimands defined in terms of $\Psi(a,a',a'',a''')$, 
under the following expanded set of ignorability assumptions: 

\begin{assumption}\label{assm:seq_ignorability_2med} 
\indent
\begin{enumerate}[noitemsep,topsep=0pt,label=(\roman*)]
	\item $Y_{i}(a,m_{1},m_{2}) \indep T_{i} | C_{i}$ for all 
	$a \in \mathcal{T}$, $m_{1} \in \mathcal{M}_{1}$, $m_{2} \in \mathcal{M}_{2}$. 
	\item $M_{1,i}(a) \indep T_{i} | C_{i}$ for all $a \in \mathcal{T}$. 
	\item $M_{2,i}(a) \indep T_{i} | C_{i}$ for all $a \in \mathcal{T}$. 
	\item $Y_{i}(a,m_{1},m_{2}) \indep M_{2,i} | C_{i}$ 
	for all $a \in \mathcal{T}$, $m_{1} \in \mathcal{M}_{1}$, $m_{2} \in \mathcal{M}_{2}$. 
	\item $Y_{i}(a,m_{1},m_{2}) \indep M_{1,i} | C_{i}$ 
	for all $a \in \mathcal{T}$, $m_{1} \in \mathcal{M}_{1}$, $m_{2} \in \mathcal{M}_{2}$. 
	\item $M_{2,i}(a,m_{1}) \indep M_{1,i} | C_{i}$ 
	for all $a \in \mathcal{T}$, $m_{1} \in \mathcal{M}_{1}$
\end{enumerate}
where $0 < \mathbb{P}[T = a | C]$, $0 < \mathbb{P}[M_{1} = m_{1} | T = a, C]$,
$0 < \mathbb{P}[M_{2} = m_{2} | T = a, M_{1} = m_{1}, C]$
for all $a \in \mathcal{T}$, $m_{1} \in \mathcal{M}_{1}, m_{2} \in \mathcal{M}_{2}$. 
\end{assumption}

Assumption~\ref{assm:seq_ignorability_2med} is a simple generalization of parts (i)-(iii) of 
Assumption~\ref{assm:seq_ignorability} to the two causally ordered mediators case. It rules out
any treatment-mediator confounding and any mediator-mediator confounding after controlling for 
pre-treatment covariates. 

For clarity, we outline an experimental procedure that can hypothetically be implemented to identify
both Lin-VanderWeele interventional effects and a number of other potentially substantively interesting
interventional effects. In particular, this experimental procedure
can identify any estimands of interest that are defined in terms of $\Psi(a,a',a'',a''')$, and the 
four terms in the Lin-VanderWeele decomposition above are a subset of these possible estimands. 

We now suppose that in addition to the ``headache drug'' defined for the previous experiment, there 
exists a similar ``insomnia drug'' that allows us to exactly set the level of insomnia a patient
experiences. We set up the following 22-arm trial (see Table~\ref{tab:rct_two_mediator}):

\begin{experiment}\label{exp:two_mediators}
\indent 
\begin{enumerate}[noitemsep,topsep=0pt]
	\item Patients in \textbf{Arm 1} receive the anti-depressant, while patients in \textbf{Arm 2} do not. 
	The levels of insomnia, headache, and depression are recorded for each patient.
	\item After recording the results of the experiment in stage 1, patients in \textbf{Arms 3-6} 
	receive combinations of treatment levels and a randomly selected level of insomnia 
	from the distribution of observed insomnia levels in either Arm 1 or Arm 2, similar to in 
	Experiment~\ref{exp:one_mediator}. The levels of insomnia, headache, and depression are recorded 
	for each patient.
	\item After recording the results of the experiment in stage 2, patients in \textbf{Arms 7-22} 
	receive combinations of treatment levels, a randomly selected level of insomnia from the distribution of
	observed insomnia levels in either Arm 1 or Arm 2, and a selected level of headache from the 
	distribution of observed headache levels in one of Arms 3-6. 
\end{enumerate}
\end{experiment}

\begin{table}[ht]
\centering
\caption{Table of results for Experiment~\ref{exp:two_mediators}}
\label{tab:rct_two_mediator}
\begin{tabularx}{\textwidth}{@{}
  >{\hsize=0.50\hsize\centering\arraybackslash}X
  >{\hsize=0.50\hsize\centering\arraybackslash}X
  >{\hsize=1.25\hsize\centering\arraybackslash}X
  >{\hsize=1.25\hsize\centering\arraybackslash}X
  >{\hsize=1.25\hsize\centering\arraybackslash}X
  >{\hsize=1.25\hsize\centering\arraybackslash}X @{}}
\hline
\addlinespace[6pt]
\makecell[bc]{\textbf{Stage}} &
\makecell[bc]{\textbf{Arm}} &
\makecell[bc]{\textbf{Treatment}} &
\makecell[bc]{\textbf{Insomnia}} &
\makecell[bc]{\textbf{Headaches}} &
\makecell[bc]{\textbf{Depression}} \\
\addlinespace[3pt]
\noalign{\hrule height 1pt}
\addlinespace[3pt]
\multirow{2}{*}{1} 
&1 & 1 & $\hat{M_{1}}(1)_{1}$ & $\hat{M_{2}}(1,\hat{M_{1}}(1)_{1})_{1}$ & $y_{1}$\\
&2 & 0 & $\hat{M_{1}}(0)_{2}$ & $\hat{M_{2}}(0,\hat{M_{1}}(0)_{2})_{2}$ & $y_{2}$\\
\addlinespace[2pt]
\noalign{\hrule height 0.5pt}
\addlinespace[2pt]
\multirow{4}{*}{2}
&3 & 1 & $\hat{G}_{1}(1)_{3}$ & $\hat{M_{2}}(1,\hat{G}_{1}(1)_{3})_{3}$ & $y_{3}$\\
&4 & 1 & $\hat{G}_{1}(0)_{4}$ & $\hat{M_{2}}(1,\hat{G}_{1}(0)_{4})_{4}$ & $y_{4}$\\
&5 & 0 & $\hat{G}_{1}(1)_{5}$ & $\hat{M_{2}}(0,\hat{G}_{1}(1)_{5})_{5}$ & $y_{5}$\\
&6 & 0 & $\hat{G}_{1}(0)_{6}$ & $\hat{M_{2}}(0,\hat{G}_{1}(0)_{6})_{6}$ & $y_{6}$\\
\addlinespace[2pt]
\noalign{\hrule height 0.5pt}
\addlinespace[2pt]
\multirow{16}{*}{3}
&7 & 1 & $\hat{G}_{1}(1)_{7}$ & $\hat{G_{2}}(1,\hat{G}_{1}(1)_{3})_{7}$ & $y_{7}$\\
&8 & 1 & $\hat{G}_{1}(1)_{8}$ & $\hat{G_{2}}(1,\hat{G}_{1}(0)_{4})_{8}$ & $y_{8}$\\
&9 & 1 & $\hat{G}_{1}(1)_{9}$ & $\hat{G_{2}}(0,\hat{G}_{1}(1)_{5})_{9}$ & $y_{9}$\\
&10 & 1 & $\hat{G}_{1}(1)_{10}$ & $\hat{G_{2}}(0,\hat{G}_{1}(0)_{6})_{10}$ & $y_{10}$\\
&11 & 1 & $\hat{G}_{1}(0)_{11}$ & $\hat{G_{2}}(1,\hat{G}_{1}(1)_{3})_{11}$ & $y_{11}$\\
&12 & 1 & $\hat{G}_{1}(0)_{12}$ & $\hat{G_{2}}(1,\hat{G}_{1}(0)_{4})_{12}$ & $y_{12}$\\
&13 & 1 & $\hat{G}_{1}(0)_{13}$ & $\hat{G_{2}}(0,\hat{G}_{1}(1)_{5})_{13}$ & $y_{13}$\\
&14 & 1 & $\hat{G}_{1}(0)_{14}$ & $\hat{G_{2}}(0,\hat{G}_{1}(0)_{6})_{14}$ & $y_{14}$\\
&15 & 0 & $\hat{G}_{1}(1)_{15}$ & $\hat{G_{2}}(1,\hat{G}_{1}(1)_{3})_{15}$ & $y_{15}$\\
&16 & 0 & $\hat{G}_{1}(1)_{16}$ & $\hat{G_{2}}(1,\hat{G}_{1}(0)_{4})_{16}$ & $y_{16}$\\
&17 & 0 & $\hat{G}_{1}(1)_{17}$ & $\hat{G_{2}}(0,\hat{G}_{1}(1)_{5})_{17}$ & $y_{17}$\\
&18 & 0 & $\hat{G}_{1}(1)_{18}$ & $\hat{G_{2}}(0,\hat{G}_{1}(0)_{6})_{18}$ & $y_{18}$\\
&19 & 0 & $\hat{G}_{1}(0)_{19}$ & $\hat{G_{2}}(1,\hat{G}_{1}(1)_{3})_{19}$ & $y_{19}$\\
&20 & 0 & $\hat{G}_{1}(0)_{20}$ & $\hat{G_{2}}(1,\hat{G}_{1}(0)_{4})_{20}$ & $y_{20}$\\
&21 & 0 & $\hat{G}_{1}(0)_{21}$ & $\hat{G_{2}}(0,\hat{G}_{1}(1)_{5})_{21}$ & $y_{21}$\\
&22 & 0 & $\hat{G}_{1}(0)_{22}$ & $\hat{G_{2}}(0,\hat{G}_{1}(0)_{6})_{22}$ & $y_{22}$\\
\addlinespace[3pt]
\noalign{\hrule height 1pt}
\end{tabularx}

\vspace{4pt}
{\footnotesize
\emergencystretch=3em
\setlength{\parindent}{1.5em}
\textsuperscript{1} $\hat{M}_{1}(1)$ and $\hat{M}_{1}(0)$ represent the distributions of insomnia in
arms 1 and 2, respectively. $\hat{G}_{1}(1)$  and $\hat{G}_{1}(0)$ represent the distributions of assigned
insomnia from random draws from $\hat{M}_{1}(1)$ and $\hat{M}_{1}(0)$, respectively. $\hat{M}_{2}(1,\hat{G}_{1}(1))$,
 $\hat{M}_{2}(1,\hat{G}_{1}(0))$,
$\hat{M}_{2}(0,\hat{G}_{1}(1))$, and $\hat{M}_{2}(0,\hat{G}_{0}(1))$ represent the distributions 
of headache in arms 3-6. $\hat{G}_{2}(1,\hat{G}_{1}(1))$, $\hat{G}_{2}(1,\hat{G}_{1}(0))$,
$\hat{G}_{2}(0,\hat{G}_{1}(1))$, and $\hat{G}_{2}(0,\hat{G}_{0}(1))$ represent the distributions of assigned
headache from random draws from the corresponding empirical distributions of $M_{2}$. Finally,
$y$ gives the observed average level of depression in each arm.\par}
\end{table}

Such an experiment fulfills all the requirements of Assumption~\ref{assm:seq_ignorability_2med}, so any
estimand written in terms of $\Psi(a,a',a'',a''')$ can be identified. This includes the four interventional 
effect analogues of the path-specific effects defined by~\cite{linvanderweele2017}. If we were only interested in 
these four interventional effects, we could reduce the number of arms in the trial and compute the 
following quantities: 
\begin{itemize}[noitemsep,topsep=0pt]
	\item $T \rightarrow Y$: $\Psi(1,0,0,0) - \Psi(0,0,0,0) = y_{14} - y_{22}$
	\item $T \rightarrow M_{1} \rightarrow Y$: $\Psi(1,1,0,0) - \Psi(1,0,0,0) = y_{10} - y_{14}$
	\item $T \rightarrow M_{2} \rightarrow Y$: $\Psi(1,1,1,0) - \Psi(1,1,0,0) = y_{8} - y_{10}$
	\item $T \rightarrow M_{1} \rightarrow M_{2} \rightarrow Y$: $\Psi(1,1,1,1) - \Psi(1,1,1,0) = y_{7} - y_{8}$
\end{itemize}

However, there are many 
other estimands defined in terms of comparisons of two values of $\Psi(a,a',a'',a''')$ that we might
be interested in, and which estimands we are interested in should
be driven by substantive considerations. As an example, consider the interventional effect given by 
$\Psi(1,0,0,0) - \Psi(0,0,0,0)$. This corresponds to the difference in the effect of a three-part intervention
to assign treatment to a patient while drawing their level of insomnia from $M_{1}(0)$ and their level of 
headache from $M_{2}(0,0)$, compared to the same intervention but with no treatment assigned. There is no 
particular reason why this particular treatment regime comparison should be of interest compared to the 
same comparison but with insomnia and headache being drawn from $M_{1}(1)$ and $M_{2}(1,1)$, respectively. 


\subsection{Vansteelandt-Daniel Interventional Effects}

~\cite{vansteelandtdaniel2017} proposes defining interventional effects for multiple mediators in
a way that both does not require knowing the causal ordering of the mediators and sums to the total
effect even in the presence of treatment-induced confounding. The distinctive feature of these effects
is that for the indirect effects through $M_{1}$ and $M_{2}$, the mediators are treated as if they
are independent regardless of the actual dependence between them.~\cite{vansteelandtdaniel2017} 
defines four effects, which we give in the context of our anti-depressant example. However, it
is now imagined that we do not know the precise dependence between insomnia and headache: we 
believe that there might exist feedback between headache and insomnia, where the existence of 
either can cause the other at a later date. We define the following effects: 

\begin{itemize}[noitemsep,topsep=0pt]
\item The \textit{interventional direct effect} is defined to be the average effect of a 
joint intervention on the treatment and two mediators that assigns the anti-depressant to
a patient while assigning a random draw of insomnia and headache from the joint distribution 
of insomnia and headache for patients that did not take the anti-depressant, compared to an 
intervention that draws from the same mediator distribution but assigns no anti-depressant. 
\item The \textit{interventional indirect effect through $M_{1}$} is defined to be the average effect
of a joint intervention on the treatment and two mediators that assigns anti-depressant to
a patient while assigning a random draw of insomnia from the corresponding distribution for patients 
who took the anti-depressant and a random draw of headache from the distribution for patients
who did not take the anti-depressant, compared to an 
intervention that also assigns the anti-depressant and assigns a level of headache from the 
distribution of headache for patients who did not take the anti-depressant but now also 
assigns a level of insomnia from the distribution of insomnia for patients who did take the 
anti-depressant. 
\item The \textit{interventional indirect effect through $M_{2}$} is defined to be the average effect
of a joint intervention on the treatment and two mediators that assigns anti-depressant to
a patient while assigning a random draw of insomnia from the corresponding distribution for patients 
who did not take the anti-depressant and a random draw of headache from the distribution for patients
who took the anti-depressant, compared to an 
intervention that also assigns the anti-depressant and assigns a level of insomnia from the 
distribution of insomnia for patients who did not take the anti-depressant but now also 
assigns a level of headache from the distribution of headache for patients who did take the 
anti-depressant. 
\item Finally,~\cite{vansteelandtdaniel2017} defines a \textit{remainder term} that corresponds to the 
difference between the (non-interventional) total effect and the sum of the three effects defined 
above. They say this captures the indirect effect resulting from the effect of the treatment on
the dependence between the counterfactuals $M_{1}(1)$ and $M_{2}(1)$ given $C$. 
\end{itemize}

These effects are also identified under Assumption~\ref{assm:seq_ignorability_2med}. However,
there are a number of important differences between these effects and the effects defined 
by~\cite{linvanderweele2017}. Vansteelandt-Daniel indirect effects compare random interventions 
on the mediator that set 
their mediator value to a random draw from its marginal distribution under either treatment or
control. In our example, this means that we imagine interventions where headache level is drawn from 
one of two marginal distributions: either the distribution of headache among a population that 
received the anti-depressant or the distribution among those who did not. For the indirect effects we 
have defined, meanwhile, we imagine drawing a mediator from a 
distribution that is conditional on both the treatment level and any mediators that we assume to 
causally precede the mediator. Thus, we imagine an intervention where headache can be drawn from 
one of four populations, where each population either was or was not assigned the anti-depressant and 
received a random level of insomnia from either a population that did or did not 
receive the anti-depressant. 

Another important difference between these effects and Lin-VanderWeele effects is in the estimation of 
the direct effect $T \rightarrow Y$, where 
both effects compare interventions where the intervention on the mediators is held constant while
varying whether treatment or control is taken, but the exact intervention on the mediators is 
different. In particular, the interventional effects that map to the Lin-VanderWeele path-specific effect 
$T \rightarrow Y$ involve drawing each mediator independently, where the first mediator is drawn
from a distribution conditional on no treatment and the second mediator is drawn from a distribution
conditional on no treatment and the first mediator having been drawn again from the distribution
conditional on no treatment (this is a separate draw from this distribution from the draw
for the value of the first mediator). The Vansteelandt-Daniel direct effect, meanwhile, involves 
only one random draw from the joint distribution of the mediators conditional on no treatment. 
This is different from both the procedure for the Lin-VanderWeele direct effect and the 
Vansteelandt-Daniel indirect effects. 

Finally, while both Vansteelandt-Daniel and Lin-VanderWeele effects are motivated by a decomposition of a total 
effect, Lin-VanderWeele effects
decompose the total interventional effect $\Psi(1,1,1,1) - \Psi(0,0,0,0)$, which is not always equal to the total
causal effect, while Vansteelandt-Daniel effects 
decompose the total causal effect $\mathbb{E}[Y_{i}(1) - Y_{i}(0)]$. As was the case for Lin-VanderWeele effects,
alternative decompositions can be pursued, and the choice between which particular decomposition to favor 
is difficult to justify. In general, we see this difference as less important given that the decomposition of 
the total effect in terms of interventional effects, while seemingly a motivation for the definition of both 
Lin-VanderWeele and Vansteelandt-Daniel effects, does not map onto meaningful information about causal 
mechanisms (see Section~\ref{subsection:interpretation}). 

We can identify Vansteelandt-Daniel effects using a version of the 
simpler two-stage experimental setup discussed in~\cite{morenobetancur2018}:

\begin{experiment}\label{exp:vd_effects}
\indent 
\begin{enumerate}[noitemsep,topsep=0pt]
	\item Patients in \textbf{Arm 1} receive the anti-depressant, while patients in \textbf{Arm 2} do not. 
	The levels of insomnia, headache, and depression are recorded for each patient.
	\item After recording the results of the experiment in stage 1, patients in \textbf{Arms 3-5} 
	receive treatment. In Arm 3, their insomnia level is drawn from the marginal distribution
	for patients in Arm 1 and headache is drawn from the marginal distribution from Arm 2; in Arm 4, insomnia is 
	drawn from Arm 2 and headache from Arm 1; and in Arm 5, both insomnia and headache are (independently)
	drawn from Arm 2. The levels of insomnia, headache, and depression are recorded for each patient. 
	To identify the direct effect, patients in \textbf{Arm 6} receive treatment 
	and a random draw from the joint distribution of insomnia and headache of patients in Arm 2, 
	while patients in \textbf{Arm 7} do not receive treatment and are also assigned a random draw 
	from the joint distribution of the mediators in Arm 2. 
\end{enumerate}
\end{experiment}

Indirect effect identification involves a comparison of the outcomes of Arm 3 and Arm 5 for $M_{1}$ and 
Arm 4 and Arm 5 respectively; direct effects can be identified by comparing outcomes for Arms 6
and 7; and the remainder term is identified by the difference between Arm 1 and Arm 2 minus the sum of 
the three other effects~\citep{morenobetancur2018}. 

\section{Interpretation of Interventional Effects}\label{subsection:interpretation}

In the previous sections, we briefly covered the distinction between the descriptive and prescriptive 
motivations underlying natural and controlled effects given by~\cite{pearl2001}. In 
Section~\ref{chapter:natural_effects}, we gave an example of why we might be interested in the 
descriptive information natural effects give us about the mechanisms underlying a causal effect. 

There are also clear examples of situations where the controlled direct effect might be of interest. To take our
anti-depressant example where headache is considered as the only mediator, to combat the negative indirect
effect of the anti-depressant on depression via headache, we might prescribe a particular dose of a 
strong anti-headache medication that is guaranteed to prevent headache. Thus, we would be testing the efficacy
of a joint intervention assigning both the anti-depressant and a particular level of headache (no headache)
on depression. This has clear clinical relevance: if such a joint intervention is seen to be effective, 
relative to a control group that took the anti-headache drug but not the anti-depressant,
this supports a treatment regime where both the anti-depressant and the anti-headache drug are exactly 
prescribed. 

Like the controlled effect, the interventional effect is motivated primarily by prescriptive rather than
descriptive considerations. By inspecting Equation~\ref{eq:RIE_one_mediator_g_comp} and comparing to the 
definition of the controlled direct effect, we can see that the interventional direct effect is the
average of controlled direct effects weighted by the probability distribution function of the mediator 
for units under no treatment. Thus, the interventional direct effect can be viewed as being itself a 
controlled direct effect under a stochastic intervention (as opposed to a deterministic one like the
controlled effect) that sets the mediator to a random level for untreated units~\cite{vanderweeleetal2014}.
In the example above, this would correspond to evaluating a joint intervention where we gave 
treated units both the anti-depressant and 
an anti-headache drug that ensured that their rates of headache merely mirrored the 
general population rather than prevented headache altogether (a more difficult assumption we would 
need to make here is that the anti-headache drug worked in a way such
that the final distribution of headache from those who had strong side effect from the anti-depressant is the 
same as for those who had only minor side effects). The results of comparing a control group to the 
group that received this joint intervention would have prescriptive implications: if this treatment regime 
were effective, the treatment regime would be an attractive intervention to implement. 

Since their introduction, interventional effects have been widely used as a substitute for natural effects 
among practitioners attempting to conduct descriptive mediation analysis~\citep{miles2023}.
This is likely because the identification 
assumptions required for interventional effects are far less stringent than natural effects, especially in 
settings with multiple mediators and treatment-induced confounding, which are very common in the social and
biomedical sciences. However,~\cite{miles2023} showed that unless very strong assumptions are made, 
interventional effects do not fulfill many important criteria that would allow us to use it as a reliable
conceptual substitute for natural effects. For instance, interventional effects do not always correctly 
detect when a mediational effect does not exist and can be of the opposite sign even when all unit-level
mediational effects are in the opposite direction.~\citep{miles2023} This means that interventional effects
are only useful in the prescriptive sense and cannot give us information about the relative strength of 
competing causal pathways in bringing about a total causal effect like natural effects can. 

The 4-way causal decomposition results that partially motivate both~\cite{linvanderweele2017} 
and~\cite{vansteelandtdaniel2017} have a descriptive motivation: they attempt to shed light on 
the relative strength of different mediated pathways in their contribution to a total effect. In light of 
the results in~\cite{miles2023}, these decompositions do not seem to be important virtues of these 
interventional effect definitions. While ambiguities in terms of choosing between different possible 
decompositions of the total causal effect exist in the natural effect decompositions discussed in 
Section~\ref{chapter:natural_effects}, these decompositions theoretically map onto the relative 
importance of different path-specific effects in a way that purely prescriptive interventional effects
do not.

Despite the conceptual issues with these decompositions, interventional effects with multiple mediators
could still be useful for measuring the effects of individual 
joint stochastic interventions. Each effect identified by Experiment~\ref{exp:two_mediators} might map to an
individually relevant stochastic intervention depending on the substantive question, without reference 
to the other effects. For example, consider the problem with insomnia and headache as causally ordered
mediators. The outcome $\Psi(1,0,0,0)$ would represent the effect of a sequential joint intervention 
that would first give treated units a sleeping pill to make their insomnia follow the distribution for 
untreated units and then give anti-headache medication to make their headaches follow the 
distribution for untreated units. If we were interested in learning how deleterious the induced side effects 
of insomnia are for mental health absent any interaction with the drug or with headache, we might instead 
be interested in $\Psi(0,1,0,0)$ and design a joint intervention to induce this effect. Similar situations 
of substantive interest might exist in cases where Vansteelandt-Daniel effects would be relevant. However, 
it is harder to think of practical 3-part joint interventions that correspond to these effects compared to 
hypothetical interventions that correspond to one mediator interventional effects. 

Even in situations where such stochastic intervention treatment regimes might be of interest, it is unclear
whether we would be interested in the exact effects given by~\cite{linvanderweele2017} 
and~\cite{vansteelandtdaniel2017}. For example, to the interventional effect $\Psi(1,0,0,0)$ 
above maps to the direct effect $T \rightarrow Y$ given by $\Psi(1,0,0,0) - \Psi(0,0,0,0)$
in~\cite{linvanderweele2017}. However, it is unclear why $\Psi(0,0,0,0)$, which represents a potentially costly
sequential joint intervention on the treatment and the two mediators, would be the relevant outcome to 
compare $\Psi(1,0,0,0)$ to to evaluate its effectiveness rather than $\mathbb{E}[Y_{i}(0,0,0)]$, the 
(deterministic) potential outcome for untreated units that we might be trying to improve upon. Additionally, 
even in situations where we posit the existence of multiple mediators, we might be interested
in one mediator~\cite{vanderweeleetal2014} interventional effects in situations where it would be costly,
unethical, or impossible to manipulate one of the mediators. 


\section{VanderWeele-Robinson Interventional Effects}\label{section:vr_int_eff}

% Need to split this apart: move some of the explanation of one mediator to here; then formally develop one
% and then multiple mediators; also might need to discuss experimental identification and assumptions, but
% putting it all together makes sense and will save time/confusion

In arguing for interpreting interventional effects as distinct stochastic interventions rather than
as mediational quantities,~\cite{miles2023} highlights a version of interventional effects developed 
by~\cite{vanderweelerobinson2014} as a particularly promising use case. These 
VanderWeele-Robinson effects are conceptually motivated by situations in which we are interested 
in what causal effect ``equalizing'' a mediator distribution between two groups
would have on an outcome of interest. This question is often of particular interest in the health 
disparities literature~\citep{jacksonvanderweele2018}.

Our running example to introduce these effects comes from~\cite{kuhaetal2021}, 
who considers how education mediates intergenerational social class mobility. Here, the ``treatment'' is
father's social class, the mediator is level of education, and the outcome is the unit's social
class at middle age. For expositional purposes, 
we treat child's social class for now as an ordered variable from ``higher'' to ``lower''
social class. When we extend these effects to two mediator problems, we will additionally 
consider child intelligence as measured by an IQ test as an additional mediator that causally precedes
level of education. 

In situations in which~\cite{vanderweelerobinson2014} are applied, the treatment is generally 
thought to be an unmanipulable characteristic such as race, sex, or social class. 
We then identify two different groups of interest --- in our running example, this might be people 
whose fathers were higher managers and 
people whose fathers were manual labourers (routine occupations). Our estimand of interest 
is then what the causal effect of randomly assigning the children of manual labourers
a level of education from the distribution of education for the children of higher managers
would have on the child's average social class. Typically, another effect of interest is the 
proportion of the difference in the outcome --- the difference between the social class of children
of managers and the social class of children of manual labourers --- that is eliminated
due to ``equalizing'' the distribution of the mediator in the manner described 
above~\citep{vanderweelerobinson2014}. 
% Add an example here or in the interpretation section: imagine a situation where we think that 
% the educational differences are driven entirely by one unjust policy; what if we repealed this policy
% (e.g. funding public schools through local property tax in the US)?
% Easy substantive interpretation relative to interventional effects

\subsection{One Mediator}\label{subsection:vr_effects_1m}

~\cite{vanderweelerobinson2014} originally introduced their version of interventional effects for single
mediator problems. 
We briefly modify the notation we have been using for interventional effects to adapt it for use
for VR interventional effects. Since it is imagined that units could not take on a different value of 
treatment, potential outcomes are a function of only the mediator $M$ and units do not have potential mediators. 
The fact that each unit
can only take on one possible value of treatment also means that there is no VanderWeele-Robinson 
interventional direct effect,
only an indirect one. Mathematically, the \textit{VR interventional indirect effect} is defined to be

% Should probably briefly introduce the parametric g-formula equation here in the single mediator 
% interventional effects section
\begin{equation}\label{eq:g_formula_vr_1m}
\begin{aligned}
 \mathbb{E}[Y_{i}(G(1)) - Y_{i}(G(0))|C, T = 0] = \\
	\int_{m} \mathbb{E}[Y_{i}(m)|T_{i} = 0, C_{i}]
	(\mathbb{P}[M = m|T = 1, C] - \mathbb{P}[M = m|T = 0, C])dm\\
\end{aligned}
\end{equation}

VR interventional effects target an analogue of the \textit{average treatment effect on the untreated} (ATU)
rather than the ATE like the interventional effects defined by~\cite{vanderweeleetal2014}. The 
VR interventional indirect effect and the interventional indirect effect are equivalent if both fulfill 
their respective identification assumptions without controlling for any confounders $C_{i}$. 
However, such effects
are identified under weaker conditions compared to interventional effects; in particular, there are 
no conditions on the independence of the treatment and the mediator and outcome variables. Thus, only
a version of condition (iii) of Assumption~\ref{assm:seq_ignorability} is required for 
identification~\citep{jacksonvanderweele2018}.

\begin{assumption}[Ignorability for VR Effects]\label{assm:seq_ignorability_vr_effects} 
\indent
$Y_{i}(m) \indep M_{i} | C_{i}$ for all $m \in \mathcal{M}$, 
where $0 < \mathbb{P}[M = m | T = 1, C]$
for $m \in \mathcal{M}$. 
\end{assumption}

The VR interventional indirect effect has a clear substantive interpretation: it represents the causal
effect of a hypothetical intervention to draw the distribution of the mediator from the treatment group
rather than the control group on the control group's outcomes. In our example, suppose our control group 
is manual labourers and our treatment group is high salaried professionals. If our 
outcome is the probability of children of father of these respective social classes becoming managers or 
high salaried professionals (destination class), we may observe a disparity in this probability depending 
on the father's social class (origin class). We may then be interested in how much this disparity in 
destination class would shrink if the children of manual labourers were assigned a level of education from
the children of high salaried professionals. This quantity could then give a measure of what effect
a policy to equalize distributions in education would have on increasing 
intergenerational social class mobility. 

% Experimental identification? Maybe not needed, can think about adding later; it makes the multiple
% mediator stuff needlessly complicated

\subsection{Multiple Mediators}\label{subsection:vr_effects_2m}

% Re-framing as reinterpretation of effects given in Jackson and VanderWeele (2018) in a situation where
% we have two causally ordered mediators rather than treating the first mediator as a confounder

% Highlight issues with marginal equalization of 2nd mediator: basically impossible to identify 
% observationally, since we would explicitly not be controlling for the first mediator (which is also
% a confounder)
We now extend VR interventional effects to the two mediator case. Extensions of the framework for 
one mediator in~\cite{vanderweelerobinson2014} have been considered in~\cite{jacksonvanderweele2018} 
in the context of considering different ways of accounting for confounders. This exercise is motivated
primarily by the health disparities literature. In problems like racial disparities in health outcomes,
certain confounders like clinical history are considered to be ``allowable'', while others like 
socioeconomic status or wealth are not~\citep{jackson2021}. Here, we consider similar effects to those
outlined in~\cite{jacksonvanderweele2018} but reframe the problem to one where instead of considering a 
mediator and confounder, we are considering two mediators of substantive interest. The difference between
our framework and that of~\cite{jacksonvanderweele2018} is similar to the difference between the framing of
treatment-induced confounding versus two mediator problems discussed in
Section~\ref{subsection:natural_effects_mm}. 

We now consider intelligence at birth as an 
additional mediator. We make the assumption that intelligence causally precedes education, so 
intelligence might both affect outcome social class directly and through its effect on educational 
attainment (see Figure~\ref{fig:application_DAG}). Let $T_{i}$ be the non-manipulable treatment, 
father's social class, let $M_{1,i}$ represent IQ with no 
corresponding potential mediators, let $M_{2,i}(m_{1})$ be potential mediators for 
education and let $Y_{i}(m_{1},m_{2})$ be potential outcomes for child's social class. 
For simplicity, we assume that no other confounders exist.

% Fit formatting
\begin{figure}
\centering
\begin{tikzpicture}

  % Define nodes
  \node[obs, align=center, minimum size=1.8cm, inner sep=0pt] (T) {Father's\\social class};
  \node[obs, right=1.5cm of T, yshift=2.5cm, align=center, minimum size=1.8cm, inner sep=0pt] (M1) {IQ};
  \node[obs, right=3.5cm of T, align=center, minimum size=1.8cm, inner sep=0pt]  (M2) {Education};
  \node[obs, right=2cm of M2, align=center, minimum size=1.8cm, inner sep=0pt]  (Y) {Social class};

  % Connect the nodes
  \draw[->, line width= 1] (T) --  (M1) ; %
  \draw[->, line width= 1] (T) --  (M2) ; %
  \draw[->, line width= 1] (M1) --  (M2) ; %
  \draw[->, line width= 1] (M1) --  (Y) ; %
  \draw[->, line width= 1] (T) to[out=325, in=215] (Y) ; %
  \draw[->, line width= 1] (M2) --  (Y) ; %

\end{tikzpicture}
\caption{Causal story for how IQ and education mediate social class mobility}
% Fix figure numbering
\label{fig:application_DAG}
\end{figure}

Suppose we are considering a problem where we are attempting to characterize how a disparity in odds 
of ending up as a high salaried professional between children of high salaried professionals and 
manual labourers is mediated by intelligence and education. 
Given this problem,~\cite{jacksonvanderweele2018} highlights a number of different VR interventional effects
we might have substantive interest in:

\begin{itemize}[noitemsep,topsep=0pt]
\item  The effect of equalizing the distribution of intelligence between the two social classes. This would 
consist of an intervention to randomly assign intelligence for children of manual labourers to a random draw
from the distribution from the marginal distribution of high salaried professionals. This is numerically 
equivalent to estimating the one mediator VR interventional effect with intelligence as the only mediator, 
given that the set of controlled confounders remains the same. 
\item  The effect of equalizing the marginal distribution of level of education between the two social classes. 
This would consist of an intervention to randomly assign level of education for children of manual labourers 
to a random draw from the marginal distribution of education among high salaried professionals. This is numerically 
equivalent to estimating the one mediator VR interventional effect with education as the only mediator, 
given that the set of controlled confounders remains the same. 
\item  The effect of equalizing the distribution of level of education, conditional on intelligence,
between the two social classes. This would consist of an intervention to randomly assign level of education 
for children of manual labourers to a random draw from the distribution of education among 
high salaried professionals with the same measured intelligence as the child of the manual labourer. 
This is numerically equivalent to estimating the one mediator VR interventional effect with education as the 
only mediator and holding the set of controlled confounders the same in addition to controlling for intelligence
(which would be framed as an additional confounder here). 
\item  The effect of equalizing the joint distribution of intelligence and education between the two social
classes. This would consist of an intervention to randomly assign intelligence and level of education 
for children of manual labourers to a random draw from the joint distribution of intelligence and education among 
high salaried professionals with the same measured intelligence as the child of the manual labourer. This is 
numerically equivalent, given the same set of confounders, to estimating the Vansteelandt-Daniel interventional 
direct effect, though the identification assumptions are less stringent compared to those effects. 
\end{itemize}

Each of the above four effects can be computed through minor modifications to Equation~\ref{eq:g_formula_vr_1m}
and can be identified under minor modifications to
Assumption~\ref{assm:seq_ignorability_vr_effects}~\citep{jacksonvanderweele2018}. One practical note to highlight is 
that the effect of equalizing the marginal level of education between social classes requires that no confounders of 
the relationship between level of education and the unit's social class exist. However, in a situation where we have
IQ and education as two mediators, IQ is a confounder of the education-destination social class relationship, so 
there is no way to observationally identify this effect: if we wanted to identify this effect, we would need to 
experimentally manipulate level of education, which might be infeasible. 

It is notable that these 
effects, with the exception of the joint intervention on intelligence and education, are analogues of the 
one mediator interventional effect definitions. While it is possible to define additional VanderWeele-Robinson 
style interventional effect analogues of the Lin-VanderWeele or Vansteelandt-Daniel effects that imagine separate, 
joint interventions on two mediators, these would be difficult to map to the types of substantive questions about 
health and social disparities that the VanderWeele-Robinson style effects are motivated by. For example, consider 
a hypothetical intervention to set the intelligence of a member of a social group to a random level of a different
(advantaged) group, but to additionally intervene to draw their level of education from their own social group's 
marginal distribution (this would correspond to a Vansteelandt-Daniel style interventional effect).~\cite{miles2023} 
shows that such a quantity could not be said to show how important intelligence is in bringing about the observed
disparity in outcomes between the two groups. Furthermore, it is difficult to think off situations where such a 
joint intervention would be feasible or interesting when compared to a simple intervention to equalize only 
intelligence and leave education unmanipulated. 


% Application section
% Information about the goals, variables, monte carlo approach (Imai et al, Kuha et al appendix)
% Illustrative purposes (no confounders): VR interventional 
% Discuss how for conditional M_2 interventional, this is actually a structural problem with any non-experimental setup
% Even more illustrative: Lin-Vanderweele and Vansteelandt-Daniel effects (compare joint intervention to me)
% Maybe: One mediator effects (for comparison)


% Notes/To-do:
% - Additive vs. proportional scale (throat-clearing)
% - Composition assumption (Vanderweele and Vansteelandt, 2009; Pearl, 2001)
% - Need to check binary vs. continuous mediator and outcome; consider at the very least switching to binary mediator, then
% giving general equations for MC estimation in section 5
% - Wording: treatment vs. exposure

% Independence of Y(a,m) and M

% Test for whether there is exposure-induced confounding


%%%%% Notes for revision -----
% - Cut down section on natural effects/potentially combine with potential outcomes stuff
% - Introduce parametric g-formula in single mediator interventional effects section
% - Add g-formula for Vansteelandt-Daniel effects
% - Explicitly mention g-formula as a term
% - Figure out what to do with the Psi(...) functions; consider either cutting entirely or adding to 
% single mediator section/incorporating in VD and VR effects sections
% - Emphasize difficulty of assumptions for conditional equalization outside of an experimental setup

%%%%% Notes for formatting clean-up -----

%%% Notation ---
% A vs. T, exposure vs. treatment!!!
% Subscripts on variables -- unit vs. population level r.v.

%%% Grammatical consistency ---
% Single vs. plural for papers with multiple authors
% Tense agreement: present vs. future, conditional vs. subjunctive, blah blah blah
% Non-manipulable vs. unmanipulable
% Multiple vs. two mediators
% Figure/table numbering and captions
% Fix quotation marks
% IQ vs. intelligence
% paper vs. dissertation vs. thesis vs. research

%%% Logical consistency ---
% IQ not evaluated at birth but instead at age 10
% Check Imai et al (a) vs. (b) cites